\songtitle{The Oath at Nashborough}

\setin{1790}
\tag*{2003}\tag{folk-ballad}\tag{original-source}\tag{british-louisiana}
\begin{notes}
Lyrics by ChatGPT based on a scene of \emph{British Louisiana} (2003), an alternate history by Carlos Thompson (Part 5, \emph{Transylvania}) -- the oath of loyalty to the King required at Nashborough to have land recognized. Retrieved from Wyomee Dank's Suno page (V6-MINI, 24 November 2026).
\end{notes}
\begin{style}
Late-1780s Southern frontier ballad; intimate, spare production with quiet fiddle beneath the words and short fiddle turns between stanzas; alternating tenor and baritone, first deliberate and on-beat, second conversational with lingering line endings, crooning legato phrasing; walking duple beat; one low, droning fiddle carries the melody, repeated refrain passed from voice to voice, first singer alone and unaccompanied for the final refrain.
\end{style}
\begin{lyrics}
\songpart{Verse 1 - First settler, tenor}
I came to Nashborough in the rain\\
To put my claim in ink.\\
The clerk pushed back his chair and said,\\
“First swear to serve the king.”

\songpart{Verse 2 - First settler, tenor}
I swore upon the book that day\\
And watched him warm the seal.\\
He asked what else I meant to claim.\\
I named the man who worked my field.

\annotation{Refrain - First settler}\\
Put your name down at Nashborough.\\
Put your mark beneath the Crown.\\
The corn will come up either way.\\
Who gets to stay when corn comes down?

\annotation{Short fiddle interlude}

\songpart{Verse 3 - Second settler, baritone}
I burned the brush and pulled the roots\\
Before they sent that clerk.\\
He can keep his book in Nashborough.\\
I’ve got a field to work.

\songpart{Verse 4 - Second settler, baritone}
You showed me where they’d put the seal.\\
I turned the paper round.\\
I knew the creek that marked your field.\\
They’d drawn it through my ground.

\annotation{Refrain - Second settler}\\
Put your name down at Nashborough.\\
Put your mark beneath the Crown.\\
The corn will come up either way.\\
Who gets to stay when corn comes down?

\songpart{Verse 5 - First settler, tenor, fiddle fades}
Come fall, you helped me haul my corn.\\
I kept the deed inside.\\
The man I’d named heard hoofbeats first.\\
We watched the rider pass us by.

\annotation{Final Refrain - First settler, unaccompanied}\\
Put your name down at Nashborough.\\
Put your mark beneath the Crown.\\
The corn will come up either way.\\
Who gets to stay when corn comes down?
\end{lyrics}

